\problemname{Treasure Hunt}

You are searching for a long-lost treasure on a massive $N \times N$ field of
cells with a magical compass. There are a total of $K \leq 3$ treasure chests
hidden at various distinct cells in the field. Your goal is simple: to find them
all!

When you place the compass on one of the cells, it will find the shortest paths
to a treasure chest using only moves left, up, right, and down (i.e. a closest
treasure chest according to Manhattan distance). It will then show the direction
of the first move. If multiple valid first moves lead to a closest treasure chest
(or several of them), the compass will return all of them. If the cell contains
treasure, the compass will indicate that instead.

Whenever you find a treasure chest, you empty the contents, but you cannot remove
the chest --- it is too heavy. Your compass does not know whether a chest is full
or empty. It will point the way to a closest chest, empty or full.

Try to find the location of all treasure chests in the fewest compass queries
possible.

\section*{Task}
This is an interactive task. In each test case (i.e. each execution of your
program), your program will have to solve several treasure hunts. You should
interact with the grader by using a library provided by the organisers. This
library contains the following declarations:

\begin{itemize}
  \item \texttt{void NextHunt(int \&N, int \&K)} --- call this function to start
        the next treasure hunt. It will set the grid size in the variable
        \texttt{N} and the number of treasures in \texttt{K}. If there are no
        more treasure hunts to be solved in the current run, the function will
        set \texttt{N} and \texttt{K} to $-1$; in that case, you should then
        terminate your program with the exit code $0$. Note that you may call
        this function before finding all the treasures on the current hunt, e.g.
        if you are trying to only solve for partial points.
  \item \texttt{enum \{ TREASURE = 0, DIR\_RIGHT = 1, DIR\_UP = 2, DIR\_LEFT = 4,
        DIR\_DOWN = 8 \};} --- these are constants used in the return values of
        the \texttt{Query} function (see below).
  \item \texttt{int Query(int x, int y)} --- if the cell at coordinates $(x, y)$
        contains a treasure, this function returns \texttt{TREASURE}. Otherwise
        it returns the sum of one or more of the constants \texttt{DIR\_RIGHT},
        \texttt{DIR\_UP}, \texttt{DIR\_LEFT}, and \texttt{DIR\_DOWN}, indicating
        which moves the compass returns when placed on the cell $(x, y)$. The
        coordinates $x$ and $y$ must be integers from $0$ to $N-1$. (Note: in
        this task, $y$-coordinates increase from top to bottom.)
\end{itemize}

After \texttt{NextHunt} sets $N = K = -1$, your program must not call either
\texttt{NextHunt} or \texttt{Query} again. It must also not call \texttt{Query}
before the first call to \texttt{NextHunt}. If your program fails to observe
these constraints, or if it makes more than $1000$ queries within a single
treasure hunt, or if it provides out-of-range values of $x$ and/or $y$ when
calling \texttt{Query}, the library will terminate your program and return a
run-time-error (RTE) verdict for the current test case.

A treasure is considered to have been found if you queried the cell containing it
at least once. If your program calls \texttt{NextHunt} before finding all the
treasures of the current hunt, this is not treated as an error, but it will affect
your score (more on scoring below).

To access the library, your program should include the header file
\texttt{treasurehuntlib.h}, that is, write \texttt{\#include "treasurehuntlib.h"}
at the top of your program.

The header file is provided as an attachment, together with
\texttt{treasurehuntlib-public.cpp}, a simple implementation of the library that
you can use while developing your solution, and \texttt{hunt\_sample.cpp}, a small
example solution. To compile the library together with your program, simply add
its name as a parameter to the compiler, e.g.

\begin{verbatim}
g++ foo.cpp treasurehuntlib-public.cpp
\end{verbatim}

if \texttt{foo.cpp} is the name of the file containing your solution.

The implementation in \texttt{treasurehuntlib-public.cpp} supports, in addition to
the above declarations, another function called
\texttt{void InitFromFile(const char *fileName)}, which reads a list of treasure
hunts from a file and lets you play the game with those instead of generating its
own random treasure hunts. You will find more details inside
\texttt{treasurehuntlib-public.cpp} itself; the attachment \texttt{sample.in} is
such a file, holding the single treasure hunt of the example below.

On the evaluation server, a different implementation of the library will be used,
so you should not make any assumptions about how exactly the implementation
works. You may, however, assume that it does not pollute the global namespace
with any declarations other than those listed above (\texttt{NextHunt},
\texttt{Query}, and the five constants).

\section*{Input}
Your program must not read from the standard input: it is used by our
implementation of the library on the evaluation server to communicate with the
rest of the evaluation environment. Everything your program needs to know arrives
through \texttt{NextHunt} and \texttt{Query}.

\begin{itemize}
  \item $2 \leq N \leq 10^6$
  \item $1 \leq K \leq 3$
  \item Within any single execution of your program, there will be at most
        $100\,000$ treasure hunts.
  \item Within any single treasure hunt, you may make at most $1000$ queries.
  \item The grading system is not adaptive.
\end{itemize}

\section*{Output}
Your program must not write to the standard output either, for the same reason.
Nothing is expected of it beyond calling the library correctly and terminating
with exit code $0$.

\section*{Scoring}
A subtask may consist of several test cases (several executions of your program),
and each test case may consist of several treasure hunts. For the purposes of
scoring, all the treasure hunts of a given subtask are evaluated together,
regardless of how they were originally split amongst the test cases. For the
$i$-th treasure hunt, let us denote the grid size by $N_i \times N_i$, the number
of treasures by $K_i$, the number of queries made by your program by $Q_i$ and
the number of treasures found by your program by $F_i$. Let us furthermore denote
by $S$ the total number of points available for this subtask. Then the number of
points awarded to your program for this subtask is:

\begin{itemize}
  \item If your program always found all the treasures (i.e. if $F_i = K_i$ for
        all $i$), its score depends on
        $t_i = \frac{Q_i}{\lceil \log_2 N_i \rceil}$:
        \[
          \frac{S}{2} + \frac{S}{2} \cdot \min_i f(t_i) \text{ points, where }
          f(t_i) = \begin{cases}
            1, & t_i \leq 11 \\
            1 - (t_i - 11)/9, & 11 \leq t_i \leq 20 \\
            0, & t_i \geq 20.
          \end{cases}
        \]
  \item If your program did not always find all the treasures (i.e. there exists
        an $i$ such that $F_i < K_i$), it receives
        $\frac{S}{2} \cdot \min_i \frac{F_i}{K_i}$ points.
\end{itemize}

In other words, you get half of the points for finding all the treasures, and the
other half for finding them in as few queries as possible. For a perfect score,
your solution should find all the treasures in at most
$11 \lceil \log_2 N_i \rceil$ queries. Between $11 \lceil \log_2 N_i \rceil$ and
$20 \lceil \log_2 N_i \rceil$ queries, the score decreases linearly; and if your
solution makes more than $20 \lceil \log_2 N_i \rceil$ queries, it will only get
the first half of the points for finding all the treasures. (The symbols
$\lceil \cdot \rceil$ mean that the value of $\log_2 N_i$ is rounded up to the
nearest integer.)

If the formulas above result in a non-integer number of points for the subtask, it
will be rounded to the nearest integer.

If your program has a runtime error or doesn't adhere to the aforementioned
protocol in using the library, it will get $0$ points for the whole subtask. To
receive partial points for finding a subset of the treasures, your program
therefore needs to finish the search gracefully by calling \texttt{NextHunt}.

\noindent
\begin{tabular}{| l | l | p{12cm} |}
  \hline
  \textbf{Group} & \textbf{Points} & \textbf{Constraints} \\ \hline
  $1$ & $10$ & $K = 1$ \\ \hline
  $2$ & $30$ & $K = 2$ \\ \hline
  $3$ & $60$ & $K = 3$ \\ \hline
\end{tabular}

\section*{Example}
\noindent
\begin{tabular}{| l | l |}
  \hline
  \textbf{Call} & \textbf{Return value} \\ \hline
  \texttt{NextHunt(N, K)} & $N = 4$, $K = 1$ \\ \hline
  \texttt{Query(2, 0)}    & \texttt{DIR\_DOWN + DIR\_RIGHT} $= 9$ \\ \hline
  \texttt{Query(3, 1)}    & \texttt{DIR\_DOWN} \\ \hline
  \texttt{Query(3, 2)}    & \texttt{TREASURE} \\ \hline
  \texttt{NextHunt(N, K)} & $N = -1$, $K = -1$ \\ \hline
\end{tabular}

\section*{Comment}
There is no sample input or output: your program never reads or writes anything
itself, so the table above is the whole interaction. The single chest sits at
$(3, 2)$, and the attachment \texttt{sample.in} replays exactly this hunt against
the provided local implementation of the library.
